\documentclass[fleqn]{article}
\oddsidemargin 0.0in
\textwidth 6.0in
\thispagestyle{empty}
\usepackage{import}
\usepackage{amsmath}
\usepackage{graphicx}
\usepackage{flexisym}
\usepackage{amssymb}
\usepackage{bigints} 
\usepackage[english]{babel}
\usepackage[utf8x]{inputenc}
\usepackage{float}
\usepackage[colorinlistoftodos]{todonotes}

\definecolor{hwColor}{HTML}{AD53BA}

\begin{document}

  \begin{titlepage}

    \newcommand{\HRule}{\rule{\linewidth}{0.5mm}}

    \center


    \textsc{\LARGE Arizona State University}\\[1.5cm]

    \textsc{\LARGE Linear Algebra }\\[1.5cm]


    \begin{figure}
      \includegraphics[width=\linewidth]{asu.png}
    \end{figure}


    \HRule \\[0.4cm]
    { \huge \bfseries Quiz Four}\\[0.4cm] 
    \HRule \\[1.5cm]

    \textbf{Behnam Amiri}

    \bigbreak

    \textbf{Prof: Sergei Suslov}

    \bigbreak


    \textbf{{\large \today}\\[2cm]}

    \vfill

  \end{titlepage}

  \begin{enumerate}
    \item (10 points) Find the angle between the vectors $u=(-1, 2, -3, 4)$ and $v=(1, -3, 2, -4)$ and 
    verify the Cauchy-Schwarz inequality, $|\left\langle u, v\right\rangle| \leq ||u|| ~ ||v||$, for 
    them using the standard inner product in $\mathbb{R}^4: ~ \left\langle u,v\right\rangle =u_1 v_1+u_2 v_2+u_3 v_3+u_4 v_4$.

      \textcolor{hwColor}{
        $
          \\
          \\
          \begin{cases}
              \left\langle \overrightarrow{u}.\overrightarrow{v}\right\rangle=(-1)(1)+(2)(-3)+(-3)(2)+(4)(-4)=-29
              \\
              \\
              ||\overrightarrow{u}||^2=\left\langle \overrightarrow{u}.\overrightarrow{u}\right\rangle=(-1)^2+(2)^2+(-3)^2+(4)^2=30 \longrightarrow ||\overrightarrow{u}||=\sqrt{30}
              \\
              \\
              ||\overrightarrow{v}||^2=\left\langle \overrightarrow{v}.\overrightarrow{v}\right\rangle=(1)^2+(-3)^2+(2)^2+(-4)^2=30 \longrightarrow ||\overrightarrow{v}||=\sqrt{30}
          \end{cases}
          \\
          \\
          \\
          cos(\theta)=\dfrac{\left\langle \overrightarrow{u}.\overrightarrow{v}\right\rangle}{||\overrightarrow{u}||||\overrightarrow{v}||}
          =-\dfrac{29}{\sqrt{30} \sqrt{30}}
          =-\dfrac{29}{30}
          \\
          \\
          \therefore ~~~~ \theta=arccos(-\dfrac{29}{30}) \approx 165^{\circ}
          \\
          \\
          \\
          |\left\langle u, v\right\rangle| \leq ||u|| ~ ||v|| \Longrightarrow 29 \leq \sqrt{30} \sqrt{30} \Longrightarrow 29 \leq 30
        $
        \\
        \\
        As it is shown the Cauchy-Schwarz inequality holds. $\blacksquare$
      }

    \item (10 points) Verify that the set of three vectors $v_1=(0, 1, 0), v_2=(-\dfrac{5}{13}, 0, \dfrac{12}{13}),$ and \\
    $v_3=(\dfrac{12}{13}, 0, \dfrac{5}{13})$ forms an orthonormal basis for $\mathbb{R}^3$ with the standard Euclidean inner
    product. Express the vector $u=(1, 2, 3)$ as a linear combination of the vectors in the given set: 
    $u=c_1 v_1+c_2 v_2+c_3 v_3$ and find the coordinates $c_1, c_2$ and $c_3$.

      \textcolor{hwColor}{
        \\
        Based on the given vectors, let's create a set of column space: 
        $$
          S=\{ \left(0, 1, 0\right), \left(-\dfrac{5}{13}, 0, \dfrac{12}{13}\right), \left(\dfrac{12}{13}, 0, \dfrac{5}{13}\right) \} 
        $$
        Now let's find the inner product of these vectors: 
        \\
        \\
        $
          \begin{cases}
            \left\langle v_1, v_2\right\rangle=(0)(-\dfrac{5}{13})+(1)(0)+(0)(\dfrac{12}{13})=0
            \\
            \\
            \left\langle v_1, v_3\right\rangle=(0)(\dfrac{12}{13})+(1)(0)+(0)(\dfrac{5}{13})=0
            \\
            \\
            \left\langle v_2, v_3\right\rangle=(-\dfrac{5}{13})(\dfrac{12}{13})+(0)(0)+(\dfrac{12}{13})(\dfrac{5}{13})=0
          \end{cases}
        $
        \\
        \\
        These results show that the given vectors are linearly independent, therefore they form a basis for $\mathbb{R}^3$.
        Also, it is obvious that the given vectors are orthogonal to each other, therefore the set $S$ is an orthogonal
        set in $\mathbb{R}^3$. 
        \\
        \\
        \\
        $
          \begin{cases}
            ||v_1||=\sqrt{(0)^2+(1)^2+(0)^2}=1
            \\
            \\
            ||v_2||=\sqrt{(-\dfrac{5}{13})^2+(0)^2+(\dfrac{12}{13})^2}=1
            \\
            \\
            ||v_3||=\sqrt{(\dfrac{12}{13})^2+(0)^2+(\dfrac{5}{13})^2}=1
          \end{cases}
        $
        \\
        \\
        Therefore, the set $S$ is orthonormal.
        \\
        \\
        Time to find the coordinates $c_1, c_2,$ and $c_3$.
        \\
        \\
        $
          u=c_1 v_1+c_2 v_2+c_3 v_3 \Longrightarrow (1, 2, 3)=c_1 (0, 1, 0)+c_2 (-\dfrac{5}{13}, 0, \dfrac{12}{13})+c_3 (\dfrac{12}{13}, 0, \dfrac{5}{13})
          \\
          \\
          \\
          \Longrightarrow (1, 2, 3)=(0, c_1, 0)+(-\dfrac{5}{13}c_2, 0, \dfrac{12}{13}c_2)+(\dfrac{12}{13}c_3, 0, \dfrac{5}{13}c_3)
          \\
          \\
          \\
          \begin{cases}
            -\dfrac{5}{13}c_2+\dfrac{12}{13}c_3=1
            \\
            \\
            c_1=2
            \\
            \\
            \dfrac{12}{13}c_2+\dfrac{5}{13}c_3=3 
          \end{cases} \Longrightarrow
          \begin{cases}
            -\dfrac{5}{13}c_2+\dfrac{12}{13}c_3=1
            \\
            \\
            \dfrac{12}{13}c_2+\dfrac{5}{13}c_3=3 
          \end{cases} \Longrightarrow 
          \left(\begin{array}{cc|c}  
            -\dfrac{5}{13} & \dfrac{12}{13} & 1
            \\
            \\
            \dfrac{12}{13} & \dfrac{5}{13} & 3
          \end{array}\right)
          \\
          \\
          \rule{15cm}{3pt}
          \\
          \\
          -\dfrac{13}{5}R_1 \rightarrow R_1
          \\
          \\
          \left(\begin{array}{cc|c}  
            1 & -\dfrac{12}{5} & -\dfrac{13}{5}
            \\
            \\
            \dfrac{12}{13} & \dfrac{5}{13} & 3
          \end{array}\right)
          \\
          \\
          \rule{15cm}{3pt}
          \\
          \\
          -\dfrac{12}{13}R_1+R_2 \rightarrow R_2
          \\
          \\
          \left(\begin{array}{cc|c}  
            1 & -\dfrac{12}{5} & -\dfrac{13}{5}
            \\
            \\
            0 & \dfrac{13}{5} & \dfrac{27}{5}
          \end{array}\right)
          \\
          \\
          \rule{15cm}{3pt}
          \\
          \\
          \dfrac{5}{13}R_2 \rightarrow R_2
          \\
          \\
          \left(\begin{array}{cc|c}  
            1 & -\dfrac{12}{5} & -\dfrac{13}{5}
            \\
            \\
            0 & 1 & \dfrac{27}{13}
          \end{array}\right)
          \\
          \\
          \rule{15cm}{3pt}
          \\
          \\
          \dfrac{12}{5}R_2+R_1 \rightarrow R_1
          \\
          \\
          \left(\begin{array}{cc|c}  
            1 & 0 & \dfrac{31}{13}
            \\
            \\
            0 & 1 & \dfrac{27}{13}
          \end{array}\right)
          \\
          \\
          \\
          \therefore ~~~~~ \begin{cases}
            c_1=2
            \\
            \\
            c_2=\dfrac{31}{13}
            \\
            \\
            c_3=\dfrac{27}{13}
          \end{cases} ~~~~~ \checkmark
        $
      }

    \item (10 points) Prove the Cauchy-Schwarz inequality, $|\left\langle u, v\right\rangle| \leq ||u|| ~ ||v||$,
    for two arbitrary vectors in an abstract inner product space $V$. Write this inequality explicitly in $\mathbb{R}^4$
    with the standard inner product: $\left\langle u,v\right\rangle =u_1 v_1+u_2 v_2+u_3 v_3+u_4 v_4$.


    \item (10 points) Convert the basis $v_1=(1, -1, 0), v_2=(0, 1, -1),$ and $v_3=(-1, 1, -1)$ for $\mathbb{R}^3$ into
    an orthonormal basis, using the Gram-Schmidt process and the standard inner product in $\mathbb{R}^3$.


    \item (Extra credit, 5 points) The linear transformation, $L \left[p(x)\right]=\dfrac{d}{dx}p(x)+p(0)$, maps 
    a polynomial $p(x)$ of degree $\leq 2$  into a polynomial of degree $\leq 1$, namely, $L: P_2 \rightarrow P_1$.
    Find the matrix representation of $L$ with respect to the ordered bases $\{ x^2, x, 1\}$ and $\{ x, 1\}$.

    \item (10 points) Show that for any two vectors $u$ and $v$ in an inner product space $V$,
    $$||u+v||^2+||u-v||=2\left(||u||^2+||v||^2\right)$$
    Give a geometric interpretation of this result for the vector space $\mathbb{R}^2$.

  \end{enumerate}

\end{document}
